\begin{lemma}
Let \(0<\varepsilon\), let \(K,n,\ell_R,\ell_B\) be natural numbers with
\(0<K\), and let \(p,\eta\in\mathbb R\).  Put
\[
x_R=\frac{\ell_R}{K},\qquad x_B=\frac{\ell_B}{K},
\qquad
f=\noderef{ProjectionOpenPairFunction}.
\]
Assume
\[
0\le K\log n,\qquad \eta\le\frac{\varepsilon}{80},
\qquad 0\le\ell_R,\ell_B\le K,
\qquad x_R+x_B\le1-\frac{\varepsilon}{2}.
\]
Then
\[
\begin{aligned}
&\left(\eta+\left(\frac12+\eta\right)-\frac{2-x_R-x_B}{2}\right)K\log n\\
&\quad
+\min\left\{0,\,
-p\left(f(\ell_R,\ell_B,K,p,n)-\varepsilon^3K^2\right)\right\}
\le -4\eta K\log n.
\end{aligned}
\]
\end{lemma}
\begin{proof}
The condition \(0<K\) is the natural-number boundary condition needed to make
the displayed ratios \(x_R=\ell_R/K\) and \(x_B=\ell_B/K\) ordinary real
ratios.  The feasibility assumptions \(0\le\ell_R,\ell_B\le K\) are the range
used in the parent summation; this proof only needs the resulting lower-regime
inequality \(x_R+x_B\le1-\varepsilon/2\).

The exact definition of \(f=\noderef{ProjectionOpenPairFunction}\), with
\(\binom y2_{\mathbb R}=y(y-1)/2\), is
\[
\begin{aligned}
f(\ell_R,\ell_B,K,p,n)
&=\frac{\ell_R(\ell_R-1)}2+\frac{\ell_B(\ell_B-1)}2\\
&\quad-\min\left\{
\frac{(K-\ell_R)(K-\ell_R-1)}2,\,
\frac{(pn)(pn-1)}2+
\frac{(K-\ell_R-pn)(K-\ell_R-pn-1)}2
\right\}\\
&\quad-\min\left\{
\frac{(K-\ell_B)(K-\ell_B-1)}2,\,
\frac{(pn)(pn-1)}2+
\frac{(K-\ell_B-pn)(K-\ell_B-pn-1)}2
\right\}.
\end{aligned}
\]
No lower bound on this expression is required in the lower regime.  The
summand contains
\[
\min\left\{1,\exp\left(-p(f(\ell_R,\ell_B,K,p,n)-\varepsilon^3K^2)\right)\right\}.
\]
After taking logarithms in the parent argument, this factor contributes at
most
\[
\min\left\{0,\,
-p\left(f(\ell_R,\ell_B,K,p,n)-\varepsilon^3K^2\right)\right\}\le0, \tag{1}
\]
because the minimum of \(0\) and any real number is at most \(0\).  Thus the
rounded values of \(K\) and \(p\) introduce no error term in this lower-regime
helper; the only large-\(n\) facts needed before applying the helper are the
two parent-side boundary estimates \(0<K\) and \(0\le K\log n\).

It remains to bound the coefficient of \(K\log n\).  This coefficient is
\[
\eta+\frac12+\eta-\frac{2-x_R-x_B}{2}
=2\eta+\frac{x_R+x_B-1}{2}.
\]
Since \(x_R+x_B\le1-\varepsilon/2\),
\[
2\eta+\frac{x_R+x_B-1}{2}
\le 2\eta-\frac{\varepsilon}{4}.
\]
The hypothesis \(\eta\le\varepsilon/80\) gives
\[
6\eta\le \frac{6\varepsilon}{80}
=\frac{3\varepsilon}{40}
\le \frac{\varepsilon}{4},
\]
and therefore
\[
2\eta-\frac{\varepsilon}{4}
\le -4\eta,
\]
because the last inequality is exactly
\(6\eta\le\varepsilon/4\) after adding \(4\eta+\varepsilon/4\) to both sides.
Multiplying by the nonnegative quantity \(K\log n\) gives
\[
\left(\eta+\left(\frac12+\eta\right)-\frac{2-x_R-x_B}{2}\right)K\log n
\le -4\eta K\log n. \tag{2}
\]
Adding the nonpositive term from (1) to the left side of (2) preserves the
inequality and proves the claim.
\end{proof}
